\documentclass[12pt, a4paper]{article}

% ===================================================================
% PAQUETES ESENCIALES
% ===================================================================

% Codificación y lenguaje
\usepackage[utf8]{inputenc}
\usepackage[spanish]{babel}

% Matemáticas
\usepackage{amsmath}
\usepackage{amssymb}
\usepackage{amsfonts}

% Gráficos e imágenes
\usepackage{graphicx}
\graphicspath{{../assets/images/}}

% Configuración de página
\usepackage{geometry}
\geometry{a4paper, margin=1in}

% Enlaces e hipervínculos
\usepackage{hyperref}
\hypersetup{
    colorlinks=true,
    linkcolor=blue,
    filecolor=magenta,
    urlcolor=cyan,
    pdftitle={Título del Documento},
    pdfauthor={Autor},
}

% ===================================================================
% COMANDOS PERSONALIZADOS (Opcional)
% ===================================================================

% Derivadas
\newcommand{\dx}[1]{\frac{d}{dx}\left(#1\right)}
\newcommand{\deriv}[2]{\frac{d#1}{d#2}}

% ===================================================================
% INFORMACIÓN DEL DOCUMENTO
% ===================================================================

\title{Título del Documento}
\author{Nombre del Autor}
\date{\today}

% ===================================================================
% INICIO DEL DOCUMENTO
% ===================================================================

\begin{document}

\maketitle
\tableofcontents
\newpage

% ===================================================================
% CONTENIDO
% ===================================================================

\section{Introducción}

Escribe aquí la introducción de tu documento.

\section{Desarrollo}

\subsection{Subsección 1}

Contenido de la subsección 1.

\subsection{Subsección 2}

Contenido de la subsección 2.

\section{Conclusión}

Escribe aquí las conclusiones.

% ===================================================================
% REFERENCIAS (Opcional)
% ===================================================================

% \begin{thebibliography}{9}
% \bibitem{referencia1}
% Autor, Título del libro, Editorial, Año.
% \end{thebibliography}

\end{document}
